%% kapa: exam=2 round=21 year=2010 period=3 subject=compensation_law subject_ko=감정평가 및 보상법규 seq=2 number=2 points=30
%% source: Q-Net 감정평가사 기출문제 (cst00309) artlSeq=5024183 file=extracted/1008/5024183/2010년 제21회 감정평가사 2차시험 기출문제/3교시.hwp
%% conversion: kordoc 4.24.51 -> tools/gen2.py (auto); review=proofread-sonnet; flags=none
\begin{longquestion}{【문제 2】}
뉴타운(New Town)개발이 한창인 A지역 인근에 주택을 소유한 P는 자신의 주택에 대하여 전년도 대비 현저히 상승한 개별공시지가를 확인하고 향후 부과될 관련 세금의 상승 등을 우려하여「부동산 가격공시 및 감정평가에 관한 법률」제12조에 따른 이의신청을 하였으나 기각되었다. 이에 P는 확정된 개별공시지가에 대하여 다시 행정심판을 제기하였으나 행정심판위원회는 그 청구를 받아들이지 않았다. 그 후 P는 자신이 소유한 주택에 대하여 전년도보다 높은 재산세(부동산보유세)를 부과 받게 되었다. (30점) \par
\medskip\noindent (1) P가 이의신청과 행정심판을 모두 제기한 것은 적법한지에 대하여 설명하시오. (10점) \par
\medskip\noindent (2) P가 소유 주택에 대하여 확정된 개별공시지가가 위법함을 이유로, 그 개별공시지가를 기초로 부과된 재산세에 대한 취소청구소송을 제기할 수 있는지에 대하여 논술하시오. (20점) \par
\end{longquestion}
